\clearpage
\section{Analiza trendów}

    \subsection{Wykresy predykcjne}

    \paragraph{}
            W celu przeprowadzenia analizy zgromadziłem dane techniczne 27 dronów klasy MALE. 
            Na ich podstawie opracuję wykresy kluczowych parametrów. 
            Pierwszym etapem będzie wyznaczenie zależności czasu lotu (jako najważniejszego parametru tej klasy) od roku produkcji płatowca. 
            W dalszej kolejności przeanalizuję wpływ wydłużenia skrzydła na najważniejsze osiągi, 
            obciążenie powierzchni oraz ciąg, 
            a także zbadam korelacje między innymi parametrami projektowymi. 
            Szczegółowe dane wejściowe wykorzystane do wizualizacji znajdują się w Załączniku nr 1.

        \begin{figure}[!h]
            \centering \includegraphics[width=1\linewidth]{Wykresy/Czas_lotu_w_zależności_od_roku_produkcji.png}
            \caption{Wykres zależności czasu od roku produkcji}
        \end{figure}

        \begin{figure}[!h]
            \centering \includegraphics[width=1\linewidth]{Wykresy/Czas_lotu_w_zależności_od_obciążenia_mocy.png}
            \caption{Wykres zależności obciążenia mocy od czasu lotu}
        \end{figure}

        \begin{figure}[!h]
            \centering \includegraphics[width=1\linewidth]{Wykresy/Czas_lotu_w_zależności_od_obciążenia_powierzchni.png}
            \caption{Wykres zależności obciążenia powierzchni od czasu lotu}
        \end{figure}

        \begin{figure}[!h]
            \centering \includegraphics[width=1\linewidth]{Wykresy/Czas_lotu_w_zależności_od_wydłużenia.png}
            \caption{Wykres zależności obciążenia powierzchni od czasu lotu}
        \end{figure}

        \begin{figure}[!h]
            \centering \includegraphics[width=1\linewidth]{Wykresy/Masa_własna_startowa_w_zależności_od_masy_startowej.png}
            \caption{Wykres zależności wspołczunika masy własnej do startowej w zależnosći od masy startowej}
        \end{figure}
    
    \newpage
    \subsection{Podsumowanie predykcji}

        \paragraph{}
            Wykresy pozwoliły mi na oszacowanie minimalnej wartości najważniejszego parametru statku powietrznego
            czasu lotu a także pozostałych współczynników opisujących charakterystykę dynamiczną BSP. 
            Na podstawie wykresu przedstawiającego stosunek masy własnej do masy startowej w funkcji masy startowej wyznaczyłem współczynniki $A$ i $C$, 
            szerzej opisane w literaturze \cite{galinski}. 
            Poprawność wyznaczonych współczynników zweryfikowałem następnie poprzez porównanie ich 
            z wartościami dostępnymi w literaturze przedmiotu \cite{raymer}.

                  \begin{longtable}{| c | m{0,5\linewidth} | c |}
            \caption{Współczyniki opisujące charakterystykę dynamiczną}
            \label{table:współczynników opisujących charakterystykę dynamiczną} \\
            \hline
            \textbf{Nazwa parametru} & \multicolumn{1}{c|}{\textbf{Wartość}} & \textbf{Jednostka} \\ \hline \hline
            \endfirsthead
            
            \hline
            \endfoot
            
            Czas lotu & 42,135 & h \\ \hline
            Wydłużenie pałata & 15,468 & N/A \\ \hline
            Obciązenie powierzchni & 97,177 & kg/m\textsuperscript{2} \\ \hline
            Obciązenie mocy & 11,648 & kg/kW \\ \hline
            Współczynik A & 1,510 & N/A \\ \hline
            Współczynik C & -0,135 & N/A \\ \hline

        \end{longtable}